\subsection{Machine Learning for Intrusion Detection}

Machine learning provides a data-driven approach to intrusion detection, where models learn to distinguish between benign and malicious traffic from labeled examples~\cite{ref_ml_ids_review}. We review the key ML paradigms and algorithms relevant to this work.

\subsubsection{Supervised Learning for Classification.}
In the IDS context, supervised learning trains a classifier $f: \mathbb{R}^d \to \{c_1, c_2, \ldots, c_K\}$ on a labeled dataset of network flows, where each flow is represented by $d$ features and assigned one of $K$ classes (e.g., \texttt{BENIGN}, \texttt{DDoS}, \texttt{BruteForce}).

\paragraph{Decision Trees.}
A decision tree recursively partitions the feature space by selecting the feature and threshold that maximizes information gain (or minimizes Gini impurity) at each internal node. For a node $t$ with class distribution $p_k$ ($k = 1, \ldots, K$), the Gini impurity is:
\begin{equation}
    \text{Gini}(t) = 1 - \sum_{k=1}^{K} p_k^2
\end{equation}
Decision trees are interpretable but prone to overfitting on complex datasets.

\paragraph{Random Forest~\cite{ref_rf_breiman}.}
Random Forest is an ensemble method that constructs $B$ decision trees, each trained on a bootstrap sample of the data with a random subset of features at each split. The final prediction is determined by majority vote:
\begin{equation}
    \hat{y} = \text{mode}\{h_b(\mathbf{x}) : b = 1, \ldots, B\}
\end{equation}
where $h_b$ is the $b$-th tree. This bagging approach reduces variance and improves generalization. Random Forest also provides a natural measure of \textbf{feature importance} based on the mean decrease in impurity across all trees.

\paragraph{XGBoost~\cite{ref_xgboost}.}
XGBoost (eXtreme Gradient Boosting) builds an ensemble of trees sequentially. At each step $t$, a new tree $f_t$ is added to minimize the regularized objective:
\begin{equation}
    \mathcal{L}^{(t)} = \sum_{i=1}^{n} l(y_i, \hat{y}_i^{(t-1)} + f_t(\mathbf{x}_i)) + \Omega(f_t)
\end{equation}
where $l$ is the loss function, $\hat{y}_i^{(t-1)}$ is the prediction from the first $t-1$ trees, and $\Omega(f_t) = \gamma T + \frac{1}{2}\lambda \|w\|^2$ is a regularization term penalizing tree complexity ($T$ = number of leaves, $w$ = leaf weights). XGBoost uses second-order Taylor expansion of the loss for efficient split finding and incorporates column subsampling and shrinkage for regularization.

\subsubsection{Deep Learning.}

\paragraph{Deep Neural Networks (DNN)~\cite{ref_deep_learning}.}
A DNN consists of an input layer, multiple hidden layers, and an output layer. Each hidden layer applies a linear transformation followed by a non-linear activation function. For layer $l$:
\begin{equation}
    \mathbf{h}^{(l)} = \sigma(\mathbf{W}^{(l)} \mathbf{h}^{(l-1)} + \mathbf{b}^{(l)})
\end{equation}
where $\mathbf{W}^{(l)}$ and $\mathbf{b}^{(l)}$ are learnable weights and biases, and $\sigma$ is an activation function. Common choices include:
\begin{itemize}
    \item \textbf{ReLU}~\cite{ref_relu}: $\sigma(z) = \max(0, z)$ --- computationally efficient, mitigates vanishing gradients.
    \item \textbf{Softmax} (output layer): $\sigma(z_k) = \frac{e^{z_k}}{\sum_{j} e^{z_j}}$ --- produces a probability distribution over $K$ classes.
\end{itemize}

Training minimizes the categorical cross-entropy loss:
\begin{equation}
    \mathcal{L} = -\sum_{i=1}^{n} \sum_{k=1}^{K} y_{ik} \log(\hat{p}_{ik})
\end{equation}
using gradient-based optimization (e.g., Adam~\cite{ref_adam}). Regularization techniques such as \textbf{Dropout}~\cite{ref_dropout} (randomly zeroing neuron outputs during training) and \textbf{Batch Normalization} are used to prevent overfitting.

DNNs can learn hierarchical feature representations automatically, making them well-suited for high-dimensional network traffic data where manual feature engineering may miss complex patterns.
